\documentclass[10pt,a4paper]{amsart}
\usepackage[margin=19mm]{geometry}
\usepackage{amsmath,amssymb,mathtools}
\usepackage[scr=rsfs,cal=euler]{mathalfa}
\usepackage{xfrac}
\usepackage[unicode,hidelinks,bookmarks=false]{hyperref}
\newcommand{\ie}{i.e.\ }
\newcommand{\cf}{cf.\ }
\newcommand{\resp}{respectively\ }
\newcommand{\Subsection}{\subsection}
\providecommand{\nobreakdash}{\nobreak\hbox{-}}
\setlength{\emergencystretch}{2em}
\multlinegap0pt

\theoremstyle{plain}
\newtheorem{theorem}{Theorem}
\newtheorem{proposition}[theorem]{Proposition}
\newtheorem{lemma}[theorem]{Lemma}
\newtheorem{corollary}[theorem]{Corollary}
\theoremstyle{definition}
\newtheorem{definition}[theorem]{Definition}
\newtheorem{assumption}[theorem]{Assumption}
\theoremstyle{remark}
\newtheorem{remark}[theorem]{Remark}
\title[Closed discrete bases in countable-dimensional topological vector spaces]{Closed discrete bases in countable-dimensional topological vector spaces: Lemma 2}
\author{Ol'ga Sipacheva}
\date{Source: arXiv:2604.07948v3 (CC BY 4.0)}
\begin{document}
\maketitle

\noindent\emph{Source and licence.} The delivered work is the article shown in the title above, version arXiv:2604.07948v3 (CC BY 4.0), distributed under the Creative Commons Attribution 4.0 International licence, as recorded in the arXiv licence metadata for this exact version and shown on the abstract page saved with this item. The full licence notice, which carries the licence URL and the address of that abstract page, is the bundled file \texttt{licenses/arxiv-2604.07948-CC-BY-4.0.txt}.\par\smallskip

\noindent\textit{Editorial scope.} Counted unit: the article's Lemma 2 (source0.tex lines 241--249, source label \texttt{l2}); statement and proof are the article's own text line for line, in the article's own notation ($\lVert\boldsymbol\cdot\rVert$, $B'$, $\widehat L$, $\mathbb F_q$), with no line omitted, substituted or re-flowed. Editorial formatting only: an AMS \texttt{amsart} preamble declaring the theorem environments the retained lines use; the article is written in plain \texttt{article} style and ships no class or style file, so none is shipped or input here; sequential numbering of the retained environments in order of appearance, whereas the article prints this lemma as Lemma 2; equation numbering is not reproduced and every cross-reference inside the retained text resolves to the wrapper's numbering. Two bibliography entries are delivered, the keys transcribed character for character from the article's own \texttt{thebibliography}; the article's printed numeric citation labels are not reproduced, and the wrapper prints the citation keys instead. No asset-inclusion command, no drawing environment and no image asset occurs in this package.


\section{The counted unit: Lemma 2 and its complete original proof}

\begin{lemma}
\label{l2}
Given any countable-dimensional topological vector space $L$ over a finite field $\mathbb F_q$, 
there exists a continuous group norm $\lVert\boldsymbol\cdot\rVert$ on $L$ and a basis $B'$ 
of $L$ such that $B'$ is discrete in the topology generated by $\lVert\boldsymbol\cdot\rVert$ and either $B'$ 
is closed in the completion $\widehat L$ of $L$ with respect to $\lVert\boldsymbol\cdot\rVert$ or $\mathbf 0$ is its 
only limit point in~$\widehat L$. Moreover, the group $L$ is incomplete with respect 
to~$\lVert\boldsymbol\cdot\rVert$. 
\end{lemma}

\begin{proof}
Let $B$ be any basis of $L$; we number its elements as $b_1,b_2, \dots$\,. Any countable 
topological space, in particular, $L$, has a countable network. According to~\cite[Lemma~8.1.12]{AT} 
and~\cite{Graev50}, there exists a continuous group norm $\lVert\boldsymbol\cdot\rVert$ on~$L$, which 
generates a group topology $\tau$ on~$L$. 

Consider the set $B'=\{b'_1, b'_2, \dots\}$ defined recursively as follows: 
\begin{enumerate} 
\item[(i)] 
$b'_1= b_1$; 

\item[(ii)] 
if $n\in \mathbb N$ and $b'_1, b'_2, \dots, b'_{n}$ are already defined, then $b'_{n+1}$ is a linear 
combination $\lambda_1 b'_1 +\lambda_2 b'_2 + \dots+ \lambda_n b'_{n}+\lambda_{n+1} b_{n+1}$, where 
$\lambda_{n+1}\ne 0$, of minimum norm (if there are several linear combinations of minimum norm, then we take any 
of them). Such a $b'_{n+1}$ exists because the number of linear combinations of elements of a finite set over a 
finite field is finite.
\end{enumerate}

It is easy to see that $B'$ has the following properties.
\begin{enumerate} 
\item[(iii)]
For each $n$,  $\langle \{b'_1, \dots, b'_n\}\rangle = \langle \{b_1, \dots, b_n\}\rangle $.
\item[(iv)]
The set $B'$ is a basis in $L$. 
\item[(v)]
If $n, i_1,\dots, i_n\in \mathbb N$ and $i_1<\dots <i_n$, then $\|b'_{i_n}\|\le 
\|\lambda_1 b'_{i_1}+\dots +\lambda_n b'_{i_n}\|$ for any $\lambda_1, \dots, \lambda_n\in 
\mathbb F_q\setminus\{0\}$. 
\end{enumerate} 

Property (i) is proved by simple induction. For $n=1$, we have $\{b'_n\}=\{b_n\}$ and hence 
$\langle \{b'_n\}\rangle =\langle \{b_n\}\rangle$. Suppose that $n>1$ and 
$\langle \{b'_1, \dots, b'_k\}\rangle = \langle \{b_1, \dots, b_k\}\rangle$ for all $k<n$. 
By definition 
$$
b'_n=\lambda_1 b'_1+ \lambda_2b'_2+ \dots +\lambda_{n-1}b'_{n-1}+\lambda_n b_n,
$$ 
where $\lambda_n\ne 0$. By the induction hypothesis 
$b'_n\in \langle \{b_1, \dots, b_n\}\rangle$ and hence $\langle \{b'_1, \dots, b'_n\}\rangle \subset 
\langle \{b_1, \dots, b_n\}\rangle$. On the other hand, 
$$
-\lambda_n b_n= \lambda_1 b'_1+ \lambda_2b'_2+ \dots +\lambda_{n-1}b'_{n-1}-b'_n
$$ 
and therefore  
$b_n\in \langle \{b'_1, \dots, b'_{n}\}\rangle$, whence 
$\langle \{b_1, \dots, b_n\}\rangle \subset \langle \{b'_1, \dots, b'_n\}\rangle$. This proves~(i).



To prove (ii), it suffices to note that the decomposition of every $b'_n$ in $B$  
contains $b_n$ with nonzero coefficient and does not contain $b_k$ with $k>n$, because 
$$
b'_n=\lambda_1 b'_1+ \lambda_2b'_2+ \dots +\lambda_{n-1}b'_{n-1}+\lambda_n b_n
$$ 
and 
$b'_1, b'_2, \dots, b'_{n-1}\in \langle \{b_1, \dots, b_{n-1}\}\rangle$ by (i). Thus, no vector $b'_n\in B'$ 
can be represented as a linear combination of vectors $b'_k$ with $k<n$, which implies the  
linear independence of~$B'$.

Let us prove~(iii). Suppose that $n, i_1,\dots, i_n\in \mathbb N$, $i_1<\dots <i_n$, and 
$\lambda_1,\dots,\lambda_n\in \mathbb F_q$. By the definition of $b'_n$ we have 
$$
b'_n=\mu_1 b'_1+ \mu_2b'_2+ \dots +\mu_{n-1}b'_{n-1}+\mu_n b_n 
$$
for some $\mu_1, \dots, \mu_{n-1}\in \mathbb F_q$ and $\mu_n\in \mathbb F_q\setminus \{0\}$; therefore, 
$$
\lambda_1 b'_1+ \lambda_2b'_2+ \dots +\lambda_n b'_n = 
\nu_1 b'_1+\nu_2 b'_2 + \dots + \nu_{n-1} b'_{n-1}+ \nu_n b_n,
$$
where $\nu_i = \lambda_i + \lambda_n\cdot \mu_i$ for $i\le n-1$ and $\nu_n=\lambda_n\cdot \mu_n\ne 0$. 
 By definition the norm of $b'_n$ is no greater than the norm of any linear combination of 
$b'_1, \dots, b'_{n-1}$ and $b_n$ in which the coefficient of $b_n$ is nonzero; this 
implies the desired inequality.  

\smallskip

It follows from property (iii) that any sequence $(b'_{n_k})_{k\in \mathbb N}$ of pairwise distinct elements 
of $B'$ which is Cauchy with 
respect to the group norm $\lVert\boldsymbol\cdot\rVert$ converges to~$\mathbf 0$ in the norm topology $\tau$. 
Indeed, by the definition of a 
Cauchy sequence, for any $\varepsilon >0$, there exists an $N\in \mathbb N$ such that 
$\|b'_{n_{k}}-b'_{n_N}\|< \varepsilon$ for any $k\ge N$. Since all integers $n_k$, $k\in \mathbb N$, 
are pairwise distinct, it follows that only finitely many of them are smaller than $n_N$ and hence there exists 
an $M\ge N$ such that $n_k>n_N$ for any $k\ge M$. 
But if $n_k > n_N$, then 
$\|b'_{n_k}\|\le \|b'_{n_k}-b'_{n_N}\|$ by property~(iii) of $B'$. Therefore, 
$\|b'_{n_k}\|<\varepsilon$ for all~$k > M$. 

Thus, in the completion $\widehat L$ of $L$ with respect to $\lVert\boldsymbol\cdot\rVert$, only $\mathbf 0$ 
can be a limit point of $B'$. 

Being countable and nondiscrete, the group $L$ is incomplete with respect to $\lVert\boldsymbol\cdot\rVert$ by 
the Baire category theorem. 
\end{proof}

\begin{thebibliography}{99}
\bibitem[AT]{AT} A.~Arhangel'skii and M.~Tkachenko, \emph{Topological Groups and Related Structures} (Atlantis Press, 2008).
\bibitem[Graev50]{Graev50} M.~I. Graev, The theory of topological groups I, \emph{Usp. Mat. Nauk} \textbf{5} (2) (1950), 3--56 (in Russian).
\end{thebibliography}
\end{document}
